\chapter{Автопилот}

Автопилот --- это программа, которая читает датчики, решает, куда лететь, и управляет моторами. На Wing это \texttt{pilot\_f}, на C10 --- ядро Melis. Исходный код у них общий, поэтому прошивка C10 с отладочной информацией служит «картой» для Wing.

\section{Слои}

Прошивка C10 (287 единиц компиляции) устроена слоями, и мы нашли те же названия в бинарнике Wing.

\begin{figure}[tbp]
\centering
\begin{tikzpicture}[
  lay/.style={draw=rule, fill=paper, rounded corners=2pt, minimum width=118mm, minimum height=9.5mm, align=center, font=\sffamily\small},
  hi/.style={lay, draw=accent, fill=accentsoft},
  ar/.style={-{Stealth[length=2.2mm]}, color=accentdark, line width=0.8pt}
]
\node[hi] (a) {\textbf{Приложение полёта}: режимы, миссия, регуляторы, отказы};
\node[lay, below=3mm of a] (b) {\textbf{Библиотеки}: PID, оценка положения и высоты, ориентация};
\node[lay, below=3mm of b] (c) {\textbf{Датчики}: IMU, барометр, оптический поток, дальномер};
\node[lay, below=3mm of c] (d) {\textbf{Слой железа}: UART, I\textsuperscript{2}C, SPI, ШИМ, драйвер моторов};
\node[lay, below=3mm of d] (e) {\textbf{Связь}: протокол awlink, радио пульта, канал к камере};
\draw[ar] (e) -- (d); \draw[ar] (d) -- (c); \draw[ar] (c) -- (b); \draw[ar] (b) -- (a);
\end{tikzpicture}
\caption{Слои автопилота. Данные идут снизу вверх, команды сверху вниз.}\label{fig:layers}
\end{figure}

\section{Поток данных}

\begin{enumerate}
  \item Датчики (инерциальный модуль, барометр, оптический поток, дальномер, камера) отдают сырые значения.
  \item Библиотеки оценивают ориентацию, высоту и положение.
  \item Слой «куда лететь» выбирает цель.
  \item Режим управления (один из нескольких) переводит цель в команду регуляторам положения, затем ориентации.
  \item Драйвер моторов отдаёт мощности на четыре мотора.
  \item Отдельно работает слой отказов и журнал.
\end{enumerate}

\section{Режимы управления}

В таблице режимов прошивки C10 (по именам файлов исходников) мы нашли:

\begin{tabularx}{\textwidth}{@{}l Y@{}}
\toprule
\hdr{Режим} & \hdr{Назначение} \\ \midrule
\texttt{stabilize}, \texttt{althold}, \texttt{poshold} & Ручной полёт, удержание высоты, удержание позиции. \\
\texttt{takeoff}, \texttt{land} & Взлёт и посадка (шаги миссии). \\
\texttt{mission} & Главный режим: выполняет шаги движения по порядку. \\
\texttt{circle}, \texttt{local360}, \texttt{jump} & Круг, оборот вокруг оси, «прыжок». \\
\texttt{track\_lines} & Следование по линии (по камере). \\
\texttt{stop} & Остановка и выключение моторов. \\ \bottomrule
\end{tabularx}

\vspace{4pt}
\pr{} Номера режимов, условия входа (например, «моторы взведены, высота известна») и переходы между ними разобраны в \texttt{docs/firmware/MISSION\_AND\_MODES.md}.

\section{Миссия}

Шоу дронов --- это миссия: список шагов с метками времени. Шаг в прошивке занимает \textbf{18 байт}.

\begin{tabularx}{\textwidth}{@{}l l Y@{}}
\toprule
\hdr{Поле} & \hdr{Размер} & \hdr{Смысл} \\ \midrule
\texttt{count} & u16 & Номер шага. \\
\texttt{time} & u16 & Момент \emph{окончания} шага, в десятых долях секунды от старта. \\
\texttt{type} & u8 & Команда: 0--14 (взлёт, перемещение, высота, круг, поворот и т.д.). \\
\texttt{data[6]} & 6 $\times$ int16 & Параметры шага (клетки, сантиметры, углы). \\ \bottomrule
\end{tabularx}

\vspace{4pt}
\pr{} Это \emph{ровно} тот кадр, который мы раньше восстановили по приложению: независимое подтверждение.

\begin{zametka}[Время держит дрона, а не гонит его]
Следующий шаг начинается, когда выполнены оба условия: цель текущего шага достигнута \emph{и} до его конца осталось не больше 0,2 с. То есть время задаёт, когда дрон \emph{отпускают} дальше, а не как быстро ему лететь. \pr{} (по ядру C10)
\end{zametka}

\section{Отказы}

По строкам автопилота видно, что он умеет реагировать на: потерю радиосвязи, потерю питания, низкий заряд, перегрев, превышение угла, столкновение, блокировку мотора, потерю положения. Для ситуации «положение потеряно» предусмотрены возврат домой и двухэтапная посадка. \pr{}

\section{Журнал и параметры}

Автопилот ведёт двоичный журнал с типами записей: \texttt{IMU}, \texttt{ATT}, \texttt{POS}, \texttt{VEL}, \texttt{BAT}, \texttt{FLOW}, \texttt{MSG}, \texttt{APRIL} (метка), \texttt{COLOR} и другие, всего около сорока. Каждая запись начинается с трёх байт заголовка. Формат записи метки: время, три целых (\texttt{cx}, \texttt{cy}, \texttt{cz}), три числа курса и номер метки с флагом. \pr{}

На нашем дроне на земле в журнале были только записи \texttt{SYS} и \texttt{MSG}: записи положения пишутся, судя по всему, только в полёте. \vyv{}

Параметры хранятся списком «имя--значение» в файлах \texttt{/netPrivate/param} и \texttt{param\_calib}. У нашего дрона в \texttt{param} стоит одна строка: \texttt{MOTOR\_\_GROUP\_NUM 1.0000}. \pr{}
